\documentclass[11pt, a4paper]{article}

\usepackage[a4paper, margin=2cm]{geometry}
\usepackage{fontspec}
\usepackage[portuguese, bidi=basic, provide=*]{babel}
\babelprovide[import, onchar=ids fonts]{portuguese}
\babelprovide[import, onchar=ids fonts]{english}
\babelfont{rm}{Noto Sans}
\usepackage{amsmath}
\usepackage{amssymb}
\usepackage{booktabs}
\usepackage{array}
\usepackage{graphicx}
\usepackage{float}
\usepackage{xcolor}
\usepackage{fancyvrb}
\usepackage{enumitem}
\setlist[itemize]{label=-}
\usepackage{hyperref}
\usepackage{listings}
\lstdefinestyle{matlab}{ 
    language=Matlab, basicstyle=\ttfamily\scriptsize,
    keywordstyle=\color{blue}\bfseries,
    commentstyle=\color{green!45!black},
    stringstyle=\color{purple},
    numbers=left,
    numberstyle=\tiny\color{gray},
    frame=single,
    breaklines=true,
    showstringspaces=false,
    tabsize=4,
    columns=fullflexible,
    keepspaces=true} 
    
\newcommand{\codigomatlab}[1]{\par%
    \IfFileExists{matlab/#1.m}% 
    {\lstinputlisting[style=matlab]{matlab/#1.m}}% 
        {\noindent\fbox{\parbox[c][2.5cm][c]{0.97\textwidth}{\centering\color{gray}% 
        Espaço reservado: \texttt{\detokenize{matlab/#1.m}}}}}}

% ---------------------------------------------------------------------------
% Validacao experimental: inclui automaticamente a figura e a saida do MATLAB
% geradas por matlab/qX_*.m (pastas figs/ e out/). Se o arquivo ainda nao
% existir, aparece um quadro reservado no lugar.
% ---------------------------------------------------------------------------
\newcommand{\figmatlab}[2]{%
    \begin{figure}[H]
        \centering
        \IfFileExists{figs/#1.png}%
            {\includegraphics[width=0.72\textwidth]{figs/#1.png}}%
            {\fbox{\parbox[c][5cm][c]{0.72\textwidth}{\centering\color{gray}%
                Espaço reservado: \texttt{\detokenize{figs/#1.png}}}}}%
        \caption{#2}
    \end{figure}
}

\newcommand{\saidamatlab}[1]{\par%
    \IfFileExists{out/#1.txt}%
        {\VerbatimInput[fontsize=\footnotesize, frame=single, label={Saída do MATLAB}]{out/#1.txt}}%
        {\noindent\fbox{\parbox[c][2.5cm][c]{0.97\textwidth}{\centering\color{gray}%
            Espaço reservado: \texttt{\detokenize{out/#1.txt}}}}}%
}

\begin{document}

\noindent
\begin{minipage}[c]{0.18\textwidth}
    \includegraphics[height=\linewidth]{misc/logo_ufsc.png}
\end{minipage}
\begin{minipage}[c]{0.79\textwidth}
    \textbf{UNIVERSIDADE FEDERAL DE SANTA CATARINA\\ CAMPUS BLUMENAU}\\[0.15cm]
    Disciplina: Sistemas a Eventos Discretos\\[0.15cm]
    Adrian Gazzani Felsky dos Anjos (21105100)
\end{minipage}

\begin{center}
    \Large{\textbf{Prova 2 — Resolução (Grafos)}}
\end{center}

% ============================================================================
% 1ª QUESTÃO - DIJKSTRA
% ============================================================================
\section*{1ª Questão: Dijkstra}

\textbf{Arestas:} $0\!\to\!2=20$, $0\!\to\!3=40$, $0\!\to\!4=30$, $1\!\to\!2=10$, $2\!\to\!4=15$, $3\!\to\!4=10$, $3\!\to\!5=80$, $4\!\to\!1=82$, $4\!\to\!5=20$.

\subsection*{Execução}
\begin{center}
\begin{tabular}{cccccccc}
\toprule
\textbf{It.} & \textbf{$u$} & \textbf{0} & \textbf{1} & \textbf{2} & \textbf{3} & \textbf{4} & \textbf{5} \\
\midrule
Início & -- & $0$ (--) & $\infty$ & $\infty$ & $\infty$ & $\infty$ & $\infty$ \\
1 & $0$ & \textbf{0 (--)} & $\infty$ & $20$ (0) & $40$ (0) & $30$ (0) & $\infty$ \\
2 & $2$ & -- & $\infty$ & \textbf{20 (0)} & $40$ (0) & $30$ (0) & $\infty$ \\
3 & $4$ & -- & $112$ (4) & -- & $40$ (0) & \textbf{30 (0)} & $50$ (4) \\
4 & $3$ & -- & $112$ (4) & -- & \textbf{40 (0)} & -- & $50$ (4) \\
5 & $5$ & -- & $112$ (4) & -- & -- & -- & \textbf{50 (4)} \\
6 & $1$ & -- & \textbf{112 (4)} & -- & -- & -- & -- \\
\bottomrule
\end{tabular}
\smallskip

\end{center}

\subsection*{Resposta}
\begin{center}
\begin{tabular}{cccc}
\toprule
\textbf{Destino} & \textbf{Custo} & \textbf{Antecessor} & \textbf{Caminho mínimo} \\
\midrule
$0$ & $0$ & -- & $0$ \\
$1$ & $112$ & $4$ & $0 \to 4 \to 1$ \\
$2$ & $20$ & $0$ & $0 \to 2$ \\
$3$ & $40$ & $0$ & $0 \to 3$ \\
$4$ & $30$ & $0$ & $0 \to 4$ \\
$5$ & $50$ & $4$ & $0 \to 4 \to 5$ \\
\bottomrule
\end{tabular}
\end{center}

\subsection*{Validação experimental (MATLAB)}
\codigomatlab{q1_dijkstra}
\figmatlab{q1_dijkstra}{Q1 — Árvore de caminhos mínimos a partir de $0$.}
\saidamatlab{q1_dijkstra}

\newpage

% ============================================================================
% 2ª QUESTÃO - BELLMAN-FORD
% ============================================================================
\section*{2ª Questão: Bellman-Ford (origem $A$)}

\textbf{Arestas:} $A\!\to\!B=2$, $A\!\to\!D=5$, $A\!\to\!F=3$, $B\!\to\!E=1$, $D\!\to\!E=1$, $F\!\to\!B=-4$, $E\!\to\!C=-3$, $E\!\to\!G=3$, $C\!\to\!B=7$, $C\!\to\!G=4$, $G\!\to\!D=-1$.

\subsection*{Execução ($|V|-1 = 6$ passadas)}
\begin{center}
\begin{tabular}{cccccccc}
\toprule
\textbf{$k$} & \textbf{A} & \textbf{B} & \textbf{C} & \textbf{D} & \textbf{E} & \textbf{F} & \textbf{G} \\
\midrule
0 & $0$ & $\infty$ & $\infty$ & $\infty$ & $\infty$ & $\infty$ & $\infty$ \\
1 & $0$ & $2$ (A) & $\infty$ & $5$ (A) & $\infty$ & $3$ (A) & $\infty$ \\
2 & $0$ & $-1$ (F) & $\infty$ & $5$ (A) & $3$ (B) & $3$ (A) & $\infty$ \\
3 & $0$ & $-1$ (F) & $0$ (E) & $5$ (A) & $0$ (B) & $3$ (A) & $6$ (E) \\
4 & $0$ & $-1$ (F) & $-3$ (E) & $5$ (A) & $0$ (B) & $3$ (A) & $3$ (E) \\
5 & $0$ & $-1$ (F) & $-3$ (E) & $2$ (G) & $0$ (B) & $3$ (A) & $1$ (C) \\
6 & $0$ & $-1$ (F) & $-3$ (E) & $0$ (G) & $0$ (B) & $3$ (A) & $1$ (C) \\
\bottomrule
\end{tabular}
\smallskip

\end{center}

\noindent\textbf{7ª passada (verificação):} nenhuma distância diminui $\Rightarrow$ não há ciclo negativo.

\subsection*{Resposta}
\begin{center}
\begin{tabular}{cccc}
\toprule
\textbf{Nó} & \textbf{Custo} & \textbf{Anterior} & \textbf{Caminho mínimo} \\
\midrule
$A$ & $0$ & -- & $A$ \\
$B$ & $-1$ & $F$ & $A \to F \to B$ \\
$C$ & $-3$ & $E$ & $A \to F \to B \to E \to C$ \\
$D$ & $0$ & $G$ & $A \to F \to B \to E \to C \to G \to D$ \\
$E$ & $0$ & $B$ & $A \to F \to B \to E$ \\
$F$ & $3$ & $A$ & $A \to F$ \\
$G$ & $1$ & $C$ & $A \to F \to B \to E \to C \to G$ \\
\bottomrule
\end{tabular}
\end{center}

\subsection*{Validação experimental}
\codigomatlab{q2_bellman_ford}
\figmatlab{q2_bellman_ford}{Q2 — Árvore de caminhos mínimos a partir de $A$.}
\saidamatlab{q2_bellman_ford}

\newpage

% ============================================================================
% 3ª QUESTÃO - PRIM
% ============================================================================
\section*{3ª Questão: Prim (início em $A$)}

\textbf{Arestas:} $AB=7$, $AD=5$, $BC=8$, $BD=9$, $BE=7$, $CE=5$, $DE=15$, $DF=6$, $EF=8$, $EG=9$, $FG=11$.

\subsection*{Execução}
\begin{center}
\begin{tabular}{ccccccccc}
\toprule
\textbf{Passo} & \textbf{$u$} & \textbf{A} & \textbf{B} & \textbf{C} & \textbf{D} & \textbf{E} & \textbf{F} & \textbf{G} \\
\midrule
Início & -- & $0$ & $\infty$ & $\infty$ & $\infty$ & $\infty$ & $\infty$ & $\infty$ \\
1 & $A$ & \textbf{0 (--)} & $7$ (A) & $\infty$ & $5$ (A) & $\infty$ & $\infty$ & $\infty$ \\
2 & $D$ & -- & $7$ (A) & $\infty$ & \textbf{5 (A)} & $15$ (D) & $6$ (D) & $\infty$ \\
3 & $F$ & -- & $7$ (A) & $\infty$ & -- & $8$ (F) & \textbf{6 (D)} & $11$ (F) \\
4 & $B$ & -- & \textbf{7 (A)} & $8$ (B) & -- & $7$ (B) & -- & $11$ (F) \\
5 & $E$ & -- & -- & $5$ (E) & -- & \textbf{7 (B)} & -- & $9$ (E) \\
6 & $C$ & -- & -- & \textbf{5 (E)} & -- & -- & -- & $9$ (E) \\
7 & $G$ & -- & -- & -- & -- & -- & -- & \textbf{9 (E)} \\
\bottomrule
\end{tabular}
\smallskip

\end{center}

\subsection*{Resposta}
\begin{center}
\begin{tabular}{ccc}
\toprule
\textbf{Ordem} & \textbf{Aresta} & \textbf{Peso} \\
\midrule
1ª & $(A, D)$ & $5$ \\
2ª & $(D, F)$ & $6$ \\
3ª & $(A, B)$ & $7$ \\
4ª & $(B, E)$ & $7$ \\
5ª & $(E, C)$ & $5$ \\
6ª & $(E, G)$ & $9$ \\
\midrule
\multicolumn{2}{r}{\textbf{Peso total da AGM:}} & \textbf{39} \\
\bottomrule
\end{tabular}
\end{center}

\subsection*{Validação experimental}
\codigomatlab{q3_prim}
\figmatlab{q3_prim}{Q3 — Árvore Geradora Mínima por Prim.}
\saidamatlab{q3_prim}

\newpage

% ============================================================================
% 4ª QUESTÃO - FORD-FULKERSON
% ============================================================================
\section*{4ª Questão: Ford-Fulkerson (fonte $0$, sumidouro $5$)}

\textbf{Capacidades:} $0\!\to\!1=16$, $0\!\to\!2=13$, $1\!\to\!2=10$, $2\!\to\!1=4$, $1\!\to\!3=12$, $3\!\to\!2=9$, $2\!\to\!4=14$, $4\!\to\!3=7$, $3\!\to\!5=20$, $4\!\to\!5=4$.

\subsection*{Caminhos aumentantes}
\begin{center}
\begin{tabular}{clcc}
\toprule
\textbf{It.} & \textbf{Caminho} & \textbf{Gargalo} & \textbf{Fluxo acumulado} \\
\midrule
1 & $0 \to 1 \to 3 \to 5$ & $\min(16,12,20) = 12$ & $12$ \\
2 & $0 \to 2 \to 4 \to 5$ & $\min(13,14,4) = 4$ & $16$ \\
3 & $0 \to 2 \to 4 \to 3 \to 5$ & $\min(9,10,7,8) = 7$ & $23$ \\
\bottomrule
\end{tabular}
\end{center}

\noindent Sem caminho aumentante restante. Corte mínimo: $S=\{0,1,2,4\}$, $T=\{3,5\}$, com $c(1,3)+c(4,3)+c(4,5) = 12+7+4 = 23$.

\subsection*{Resposta}
\begin{center}
\begin{tabular}{cccc}
\toprule
\textbf{Aresta} & \textbf{Capacidade} & \textbf{Fluxo} & \textbf{Residual} \\
\midrule
$0 \to 1$ & $16$ & $12$ & $4$ \\
$0 \to 2$ & $13$ & $11$ & $2$ \\
$1 \to 2$ & $10$ & $0$ & $10$ \\
$2 \to 1$ & $4$ & $0$ & $4$ \\
$1 \to 3$ & $12$ & $12$ & $0$ (saturada) \\
$3 \to 2$ & $9$ & $0$ & $9$ \\
$2 \to 4$ & $14$ & $11$ & $3$ \\
$4 \to 3$ & $7$ & $7$ & $0$ (saturada) \\
$3 \to 5$ & $20$ & $19$ & $1$ \\
$4 \to 5$ & $4$ & $4$ & $0$ (saturada) \\
\midrule
\multicolumn{2}{r}{\textbf{Fluxo máximo:}} & \multicolumn{2}{c}{\textbf{23}} \\
\bottomrule
\end{tabular}
\end{center}

\subsection*{Validação experimental}
\codigomatlab{q4_ford_fulkerson}
\figmatlab{q4_ford_fulkerson}{Q4 — Fluxo/capacidade em cada aresta e corte mínimo.}
\saidamatlab{q4_ford_fulkerson}

\newpage

% ============================================================================
% 5ª QUESTÃO - KRUSKAL
% ============================================================================
\section*{5ª Questão: Kruskal}

\textbf{Rótulos:} $A$ (sup. esq.), $B$ (sup. centro), $C$ (sup. dir.), $D$ (inf. esq.), $E$ (base), $F$ (inf. dir.).\\
\textbf{Arestas:} $AB=2$, $BC=3$, $AD=5$, $AE=1$, $BE=7$, $CE=3$, $CF=9$, $DE=4$, $EF=4$.

\subsection*{Execução}
\begin{center}
\begin{tabular}{ccccl}
\toprule
\textbf{Passo} & \textbf{Aresta} & \textbf{Peso} & \textbf{Decisão} & \textbf{Conjuntos} \\
\midrule
Início & -- & -- & -- & $\{A\},\{B\},\{C\},\{D\},\{E\},\{F\}$ \\
1 & $(A, E)$ & $1$ & \textbf{Aceita} & $\{A,E\},\{B\},\{C\},\{D\},\{F\}$ \\
2 & $(A, B)$ & $2$ & \textbf{Aceita} & $\{A,B,E\},\{C\},\{D\},\{F\}$ \\
3 & $(B, C)$ & $3$ & \textbf{Aceita} & $\{A,B,C,E\},\{D\},\{F\}$ \\
4 & $(C, E)$ & $3$ & Rejeitada (ciclo) & $\{A,B,C,E\},\{D\},\{F\}$ \\
5 & $(D, E)$ & $4$ & \textbf{Aceita} & $\{A,B,C,D,E\},\{F\}$ \\
6 & $(E, F)$ & $4$ & \textbf{Aceita} & $\{A,B,C,D,E,F\}$ \\
7 & $(A, D)$ & $5$ & Rejeitada (ciclo) & -- \\
8 & $(B, E)$ & $7$ & Rejeitada (ciclo) & -- \\
9 & $(C, F)$ & $9$ & Rejeitada (ciclo) & -- \\
\bottomrule
\end{tabular}
\end{center}

\subsection*{Resposta}
\begin{itemize}[noitemsep]
    \item \textbf{AGM:} $\{(A, E), (A, B), (B, C), (D, E), (E, F)\}$
    \item \textbf{Peso total:} $1 + 2 + 3 + 4 + 4 = \mathbf{14}$
\end{itemize}

\subsection*{Validação experimental}
\codigomatlab{q5_kruskal}
\figmatlab{q5_kruskal}{Q5 — Árvore Geradora Mínima por Kruskal.}
\saidamatlab{q5_kruskal}

\end{document}
